\section*{Adaptaciones a la metodología CRISP-DM}
\addcontentsline{toc}{section}{Adaptaciones a la metodología CRISP-DM}

Conservamos las cuatro tareas de la fase de Modelado de CRISP-DM y sus salidas, en el mismo orden y con el mismo nombre. Lo que se adaptó es la forma de ejecutar algunas tareas, por las características del problema (varias fotos por vaca y pocas vacas):

\begin{itemize}
  \item En lugar de elegir una sola técnica, comparamos familias de modelos contra un modelo base, en iteraciones.
  \item Desde el 12 de noviembre usamos un protocolo más estricto que el inicial. Toda elección se hace con validación cruzada interna dentro del entrenamiento y el pliegue de validación solo mide, una sola vez, al final.
  \item Cada experimento se fija antes de mirar su resultado. Si una cifra se obtuvo eligiendo con la validación a la vista, se considera optimista y no cuenta.
  \item La métrica es por vaca y se reporta como n/76. Además de la salida continua, reportamos aparte la salida redondeada a 0.25.
  \item Algunas revisiones regresaron a Preparación de Datos (segmentación del fondo y recorte de anca), algo que CRISP-DM permite.
\end{itemize}
